\textbf{Problem Understanding.} Etosha National Park (22,270 km$^2$) faces critical conservation challenges: protecting endangered species (black rhino, elephant, lion) across vast, geographically complex terrain with limited resources. Traditional uniform patrol coverage fails due to resource constraints, geographic barriers, and stochastic threat patterns.

\vspace{0.2cm}

\textbf{Our Approach.} We develop a six-layer strategic allocation framework integrating: (i) multi-dimensional protection metrics via AHP-fuzzy evaluation; (ii) graph-theoretic spatial network (50 monitoring zones); (iii) integer programming optimization (MILP core); (iv) dynamic programming for temporal allocation; (v) queueing model for staffing inference; (vi) Monte Carlo simulation (10,000 trials) for robustness; (vii) game-theoretic deterrence modeling.

\vspace{0.2cm}

\textbf{Key Results.}

\begin{table}[H]
\centering
\small
\begin{tabular}{lccc}
\toprule
\textbf{Metric} & \textbf{Uniform} & \textbf{Strategic} & \textbf{Improvement} \\
\midrule
High-Priority Species Protection & 62\% & \textbf{89\%} & +27 pp \\
Critical Area Coverage & 58\% & \textbf{94\%} & +36 pp \\
Threat Detection Rate & 55\% & \textbf{81\%} & +26 pp \\
Ranger Efficiency (km$^2$/ranger) & 180 & \textbf{320} & +78\% \\
Minimum Staffing & — & \textbf{127 rangers} & — \\
\bottomrule
\end{tabular}
\end{table}

\vspace{0.2cm}

\textbf{Key Recommendations.} (1) Priority-based deployment: 70\% resources to high-priority zones (water holes, endangered species habitats). (2) Hybrid monitoring: rangers (deterrence) + UAVs (inaccessible terrain) + camera traps (continuous surveillance). (3) Science-based staffing: 127 rangers with 3-shift rotation maintains 85\% protection for high-priority species. (4) Dynamic reallocation: adapt deployment to seasonal threat variations. (5) Universal framework: identical mathematical structure validated across Africa (Etosha), Asia (Ranthambore), South America (Pantanal).

\vspace{0.2cm}

\noindent \textbf{Control Number:} IMMC26601951
